\chapter{Host fields and averaging over auxiliary primes}\label{ch:host}

This chapter covers the abelian CM fields that the archive calls hosts, and
the theorem \texttt{supplyChi}, which supplies a host with all the analytic
properties that the explicit-formula argument needs, for every primitive
Dirichlet character $\psi$, assuming only the Galois-cohomological bound
\texttt{VWinAt}. A host is cut out of a cyclotomic field by the kernel of
$\psi$, a $2$-power level and the quadratic characters of two auxiliary
primes $q,d$. The characters of the host that are not powers of $\psi$ are
twists $\xi\chi_q^i\chi_d^j$ with $(i,j)\ne(0,0)$. The archive bounds their
prime sums at a scale polylogarithmic in $q$ by choosing $q$ and then $d$
among the majority of primes for which a mean-square estimate gives the
bound (the theorem \texttt{QH.pairEx}). The output feeds the analytic
criterion of the explicit formula and consumes the window bound
\texttt{VWinAt}; both are treated in other chapters.

Throughout, $\chi_q(a)=\left(\frac aq\right)$ and $\chi_d(a)=\left(\frac
ad\right)$ are Jacobi symbols, and a \emph{primitive character} is a pair
$(N,\chi)$ with $\chi$ a primitive Dirichlet character of conductor $N$
(the archive's \texttt{PChar}).

% ------------------------------------------------------------------
\section{Frames and hosts}

\begin{definition}[frame]\label{host:frame}
\lean{Formalism.HC.Frame, Formalism.HC.Frame.A0, Formalism.HC.Frame.n0, Formalism.HC.Frame.o, Formalism.HC.Frame.GG, Formalism.HC.Frame.cc, Formalism.HC.Frame.iod, Formalism.HC.Frame.tau1, Formalism.FrobeniusOrbits.HGrp, Formalism.FrobeniusOrbits.hcl}
For $e\ge0$ let $H_e=(\Z/2^{e+2}\Z)^{\times}/\{\pm1\}$, and let
$\operatorname{hcl}_e\colon(\Z/2^{e+2}\Z)^{\times}\to H_e$ be the quotient map.
A \emph{frame} $\phi$ consists of
\begin{itemize}
\item integers $e\ge0$ and $\mathrm{ord}$, and an abelian group structure on
  $A=\{0,\dots,\mathrm{ord}-1\}$ (so $A$ is a finite abelian group of order
  $\mathrm{ord}\ge1$);
\item a surjective homomorphism $h_0\colon A\to H_e$;
\item an element $g_0\in A$ such that $h_0(kg_0)=0$ implies $kg_0=0$ for all
  $k\in\Z$, that is $\langle g_0\rangle\cap\ker h_0=0$.
\end{itemize}
Put $A_0=\langle g_0\rangle$, $n_0=|A_0|$ and $o=\mathrm{ord}/n_0=[A:A_0]$.
The \emph{Galois model} of $\phi$ is $\mathfrak G=\{\pm1\}^2\times A$
(written multiplicatively), with the elements $c=((-1,-1),0)$,
$\iota=((1,-1),0)$ and $\tau_1(t)=((1,1),t)$ for $t\in A$.
\end{definition}

\begin{definition}[raw host]\label{host:host-raw}
\lean{Formalism.HC.HostRaw, Formalism.HC.HostRaw.condN, Formalism.FrobeniusOrbits.HostHyp}
\uses{host:frame}
Let $\phi$ be a frame. A \emph{raw host} for $\phi$ consists of
\begin{enumerate}
\item primes $q,d$ with $q\equiv3\pmod 8$, $d\equiv7\pmod{16}$ and
  $\left(\frac{-d}{q}\right)=1$ (the condition \texttt{HostHyp});
\item the condition $e\ge2$;
\item an integer $Q_S>0$ with $2^{e+2}\mid Q_S$, $\gcd(Q_S,qd)=1$ and
  $\left(\frac{-d}{l}\right)=1$ for every odd prime $l\mid Q_S$;
\item a number field $E$, Galois over $\Q$, with a group isomorphism
  $\mathrm{gal}\colon\Gal(E/\Q)\to\mathfrak G$;
\item an element $\sqrt{-d}\in E$ with square $-d$, on which $\sigma$ acts
  by the second sign of $\mathrm{gal}(\sigma)$;
\item a $\Q$-embedding $j\colon E\to\Q(\zeta_{N_H})$, where $N_H=Q_Sqd$, and
  a primitive $N_H$-th root of unity $\zeta$ there;
\item a homomorphism $\mathrm{art}\colon(\Z/N_H\Z)^{\times}\to\mathfrak G$
  such that, if $\tau\in\Gal(\Q(\zeta_{N_H})/\Q)$ sends $\zeta$ to $\zeta^u$,
  then $j(\mathrm{gal}^{-1}(\mathrm{art}(u))x)=\tau(j(x))$ for all $x\in E$;
\item $\mathrm{art}(-1)=c$;
\item the $A$-component of $\mathrm{art}(u)$ maps under $h_0$ to
  $\operatorname{hcl}_e(u\bmod 2^{e+2})$;
\item if $u\equiv1\pmod{Q_Sq}$ then $\mathrm{art}(u)\in\{1,\iota\}$, and if
  $u\equiv1\pmod{Q_Sd}$ then $\mathrm{art}(u)^2=1$.
\end{enumerate}
The integer $N_H=Q_Sqd$ is the \emph{conductor} of the host.
\end{definition}

\begin{definition}[derived host properties]\label{host:host-der}
\lean{Formalism.HC.HostDer, Formalism.HC.HostCore, Formalism.HC.DerOfRaw}
\uses{host:host-raw}
For a raw host $H$, the proposition $\mathrm{HostDer}(H)$ is the conjunction
of the following, among others: $\mathrm{art}$ is surjective; $E$ is a CM
field and $\mathrm{gal}^{-1}(c)$ is complex conjugation under every complex
embedding; $[E:\Q]=4\,\mathrm{ord}$ and the maximal real subfield $F$ has
$[F:\Q]=2\,\mathrm{ord}$ and is the fixed field of $\mathrm{gal}^{-1}(c)$;
the second sign of $\mathrm{art}(u)$ is $\left(\frac ud\right)$; the set of
embeddings sending $\sqrt{-d}$ to $i\sqrt d$ is a CM type of size
$2\,\mathrm{ord}$, stable exactly under $\mathrm{gal}^{-1}(\ker\chi_D)$, with
$\chi_D$ the second sign; for every prime $p\nmid N_H$ and every prime
$P\mid p$ of $E$ the arithmetic Frobenius exists and equals
$\mathrm{gal}^{-1}(\mathrm{art}(p))$; inertia and decomposition groups at
every prime are the images under $\mathrm{art}$ of the expected subgroups;
inertia at $d$ is $\{1,\iota\}$; descriptions of splitting in
$K=\Q(\sqrt{-d})$, of the relative different of $E/K$ as $\mathfrak b\,
\bar{\mathfrak b}$ with $\mathfrak b+\bar{\mathfrak b}=(1)$, of
$\operatorname{disc}E=(\operatorname{disc}F)^2$ and of its prime divisors;
and two-sided bounds for the root discriminant
$\mathrm{rd}_E=|\operatorname{disc}E|^{1/[E:\Q]}$,
\[
2^{e+1-2^{-e}}\sqrt d\le\mathrm{rd}_E\le\tfrac{Q_S}{2}\sqrt{qd},\qquad
4\,\mathrm{ord}\le\mathrm{rd}_E^2 .
\]
The structure \texttt{HostCore} consists of sixteen of these fields:
surjectivity, the CM and degree statements, the second sign, the
Frobenius, inertia and decomposition statements, and two statements on
complex conjugation and decomposition groups.
The proposition \texttt{DerOfRaw} states that every raw host satisfies
$\mathrm{HostDer}$.
\end{definition}

\begin{lemma}[cyclotomic basics of a host]\label{host:cyc-basics}
\lean{Formalism.HC.Der.art_surj, Formalism.HC.Der.cm, Formalism.HC.Der.conj_iff, Formalism.HC.Der.finrank_E}
\uses{host:host-raw}
Let $H$ be a raw host for $\phi$. Then $\mathrm{art}$ is surjective, $E$ is
a CM field, an automorphism $\sigma$ acts as complex conjugation under a
complex embedding if and only if $\mathrm{gal}(\sigma)=c$, and
$[E:\Q]=4\,\mathrm{ord}$.
\end{lemma}
\begin{proof}
Surjectivity: every $\sigma\in\Gal(E/\Q)$ lifts along $j$ to an
automorphism of the Galois extension $\Q(\zeta_{N_H})$, which is
$\zeta\mapsto\zeta^u$ for some unit $u$; by the compatibility of
$\mathrm{art}$ with $j$ and the injectivity of $j$,
$\mathrm{gal}(\sigma)=\mathrm{art}(u)$. Since $\mathrm{art}(-1)=c$ and
$\zeta\mapsto\zeta^{-1}$ is complex conjugation on $\Q(\zeta_{N_H})$, the
automorphism $\mathrm{gal}^{-1}(c)$ is complex conjugation under every
embedding of $E$ (extend the embedding of $E$ to the cyclotomic field).
It sends $\sqrt{-d}$ to $-\sqrt{-d}$, so no embedding is real; hence $E$
is totally complex with an automorphism inducing complex conjugation
everywhere, which is a CM field. The degree is
$|\Gal(E/\Q)|=|\mathfrak G|=4\,\mathrm{ord}$.
\end{proof}

\begin{lemma}[raw hosts have the derived properties]\label{host:raw-der}
\lean{Formalism.HC.HostRaw.der, Formalism.HC.Der.hostCore}
\uses{host:host-raw, host:host-der, host:cyc-basics}
Every raw host $H$ satisfies $\mathrm{HostDer}(H)$. In particular
\texttt{DerOfRaw} holds.
\end{lemma}
\begin{proof}
\uses{host:cyc-basics}
The proof assembles the fields one by one: the sixteen fields of
\texttt{HostCore} come from \texttt{Der.hostCore}, which collects
Lemma~\ref{host:cyc-basics} and the Frobenius, inertia and decomposition
statements obtained by restricting the corresponding statements for
$\Q(\zeta_{N_H})$ along $j$. The remaining fields (among them the CM type, splitting in
$K$, the different, the discriminant, root-discriminant bounds, and the fields \texttt{ps\_root}, \texttt{ordinary}, \texttt{eta}, \texttt{principal}, \texttt{unimodular} and \texttt{norm\_mu}) are separate
lemmas of the namespace \texttt{Formalism.HC.Der}. These live in
\arch{Formalism/SharpRestriction/RamificationFrobenius/TotientGaloisClasses/InertiaGaloisSymmetry/}
and are not expanded here.
\end{proof}

\begin{definition}[host with axioms]\label{host:host-ax}
\lean{Formalism.HC.HostAx, Formalism.HC.HostAx.ofRaw}
\uses{host:host-raw, host:host-der, host:raw-der}
A \emph{host with axioms} is a raw host $H$ together with a proof of
$\mathrm{HostDer}(H)$. By Lemma~\ref{host:raw-der} every raw host extends
to one (\texttt{HostAx.ofRaw}).
\end{definition}

\begin{remark}[the derived properties are proved]\label{host:rem-der}
The name \texttt{HostAx} suggests that the CM and ramification properties
of a host are assumed. They are not: \texttt{HostAx.ofRaw} supplies them
from Lemma~\ref{host:raw-der}, so a host with axioms carries no more
hypotheses than a raw host. The proposition \texttt{DerOfRaw} is stated
but never used elsewhere in the archive; the derivation enters only through
\texttt{HostRaw.der}.
\end{remark}

\begin{definition}[ranks]\label{host:ranks}
\lean{Formalism.BasicNotions.Ranks}
\uses{host:frame}
A \emph{rank datum} $R$ is a frame together with natural numbers $k,m,t$,
with $k,m$ even, $t$ odd, $2\le m$, $m+3\le k$, $k+m+4\le t$, $3m<t$,
$m\le2\,\mathrm{ord}\le k$ and $m^2\le k+m(t+1)$. One writes $n=\mathrm{ord}$
and $f=2n$; the archive also attaches an integer $r$ to $R$.
\end{definition}

\begin{definition}[full host]\label{host:host-full}
\lean{Formalism.BasicNotions.Host}
\uses{host:ranks, host:host-ax}
For a rank datum $R$, a \emph{host} for $R$ is a host with axioms for the
frame of $R$ together with: the identity $\psi_q(u)=\left(\frac uq\right)$
for the first sign $\psi_q$ of $\mathrm{art}(u)$; an element $\sqrt{-q}\in
E$ on which $\sigma$ acts by the first sign; for every prime $l\mid Q_S$ a
unit $u\equiv1\pmod{N_H/l}$ with $\mathrm{art}(u)\neq1$; $Q_S\le n^2$;
$o^{64}\le2^e$; $q\ne3$; $q,d>4tr+1$; a bound $b$ with
$\left(\frac{-d}{p}\right)=1$ for every odd prime $p\le b$ and $4l\le b$
for every odd prime $l\mid Q_S$; a CM-type embedding $E\to\CC$; an ideal
$\mathfrak b$ of $\mathcal O_E$ splitting the different of $E/K$; a finite
set of level primes containing two distinct primes $g_0,g_1$ and two
distinct primes $p_{\delta},p'_{\delta}\equiv1\pmod{15}$; a nonzero primitive
ideal $\mathfrak c$ of $\mathcal O_K$ supported at split primes and prime to
$2,3,5,q,d$, the level primes and $Q_S$; finite sets $TT\subset A$ and
$T_2\subset\N$, an element $u_0\in A$, an integer $L_w\ge2$ with
$4L_ws<n_0$ for $s\in T_2$ and a no-relation condition on $T_2$; and a real
number $C_{T2}$.
\end{definition}

% ------------------------------------------------------------------
\section{Characters of a host}

\begin{definition}[characters of the Galois model]\label{host:chars}
\lean{Formalism.HC.Frame.Char, Formalism.HC.Frame.TrivialOnFirstFactor, Formalism.HC.Frame.Xp, Formalism.HC.Frame.Xm, Formalism.HC.Frame.psC}
\uses{host:frame}
A \emph{character} of $\phi$ is a homomorphism $\theta\colon\mathfrak
G\to\CC^{\times}$. It is \emph{trivial on the first factor} if
$\theta(\tau_1(s))=1$ for all $s\in A_0$. Let $X_+$ (resp.\ $X_-$) be the
set of characters trivial on the first factor with $\theta(\iota)=1$
(resp.\ $\theta(\iota)=-1$). Let $\omega_q$ be the character given by the
first sign of $\{\pm1\}^2$.
\end{definition}

\begin{definition}[the Dirichlet character of a Galois character]\label{host:dir}
\lean{Formalism.HC.HostRaw.dir, Formalism.HC.HostRaw.dirAt, Formalism.HC.HostRaw.IsPrimOf, Formalism.HC.HostRaw.prim, Formalism.HC.HostRaw.IsCharE}
\uses{host:host-raw, host:chars}
% uses-lean: Formalism.HC.An.Induces, Formalism.HC.PChar
\uses{an:pchar} % resolved
Let $H$ be a raw host and $\theta$ a character of $\phi$. Then
$\theta\circ\mathrm{art}$ is a Dirichlet character modulo $N_H$, written
$\mathrm{dir}(\theta)$. A Dirichlet character $\chi$ modulo $N$ \emph{is
the primitive character of $\theta$} if $\chi$ is primitive and
$\chi(a)=\theta(\mathrm{art}(a))$ for every $a$ prime to $N_H$; the
primitive character inducing $\mathrm{dir}(\theta)$ is written
$\mathrm{prim}(\theta)$. A character $\chi$ is a \emph{character of $E$} if
it is the primitive character of some $\theta$.
\end{definition}

\begin{lemma}[uniqueness of the primitive character]\label{host:prim-unique}
\lean{Formalism.HC.HostRaw.isPrimOf_iff, Formalism.HC.HostRaw.kept_iff}
\uses{host:dir}
For a primitive character $\psi$, $\psi$ is the primitive character of
$\theta$ if and only if $\psi=\mathrm{prim}(\theta)$.
\end{lemma}
\begin{proof}
% uses-lean: Formalism.HC.An.Induces.unique, Formalism.HC.An.induces_of
Two primitive characters that agree on all integers prime to a fixed
nonzero modulus are equal (\texttt{An.Induces.unique}), and
$\mathrm{prim}(\theta)$ has the defining property by construction.
\end{proof}

\begin{definition}[kept and bounded characters]\label{host:kept}
\lean{Formalism.HC.HostRaw.Kept, Formalism.HC.HostRaw.Bounded, Formalism.HC.HostRaw.KeptInXp, Formalism.HC.HostRaw.ExceptionalCharacterKept, Formalism.HC.HostRaw.BoundedCharsPointwise, Formalism.HC.HostRaw.ChebyshevCharacterBound, Formalism.HC.HostRaw.LogDerivativeBound}
\uses{host:chars, host:dir}
% uses-lean: Formalism.HC.An.gChi, Formalism.HC.An.logDerivativeAtOne
\uses{an:gchi} % resolved
Let $H$ be a raw host and $S$ a finite set of primitive characters. Write
$g_{\chi}$ for the archive's regularised prime sum of $\chi$ and
$\lambda(\chi)=L'(1,\chi)/L(1,\chi)$.
\begin{enumerate}
\item $\theta$ is \emph{kept} if its primitive character lies in $S$.
\item $\theta$ is \emph{bounded} if $\theta\in X_+$ and is not kept, or
  $\theta\in X_-$.
\item $H$ has \emph{kept characters in $X_+$} if every $\psi\in S$ is the
  primitive character of some $\theta\in X_+$.
\item For $\psi_K\in\{\text{none}\}\cup\{\text{primitive characters}\}$, $H$
  \emph{keeps the exceptional character} if $\psi_K=\psi$ implies that
  $\psi$ is the primitive character of $\omega_q$, and $\psi_K=\text{none}$
  implies that $\omega_q$ is not kept.
\item $H$ has \emph{pointwise bounds at $x$ with exponent $\eps$} if
  $|g_{\chi}(x)|\le x^{1/2+\eps}$ for the primitive character $\chi$ of every
  bounded $\theta$.
\item $H$ has the \emph{Chebyshev bound} $(x,\eps,c)$ if
  $|g_{\chi}(y)|\le y\exp(-cx^{\eps/8})$ for every character $\chi$ of $E$
  and every $y\ge\exp(x^{\eps/3})/3$.
\item $H$ has the \emph{log-derivative bound} $(S,X,x,\eps,C)$ if
  $|\lambda(\chi)|\le x^{\eps}$ for the primitive character of the second
  sign, and $|\lambda(\chi)|\le C\log X$ for the primitive character of
  every bounded $\theta$.
\end{enumerate}
\end{definition}

\begin{definition}[host conditions]\label{host:host-conditions}
\lean{Formalism.HC.HostConditions}
\uses{host:kept}
% uses-lean: Formalism.HC.HType, Formalism.HC.HConst
For a type $T$ (carrying a kept set $S$, an exceptional slot $\psi_K$ and an
exponent $\eps$) and constants $(c,C)$, the \emph{host conditions} for a raw
host $H$ at $(X,x)$ are: kept characters in $X_+$, the Chebyshev bound
$(x,\eps,c)$, the log-derivative bound $(S,X,x,\eps,C)$, pointwise bounds at
$x$ with exponent $\eps$, and the exceptional character kept.
\end{definition}

\begin{remark}[the host conditions are not the live hypothesis]\label{host:rem-hostconditions}
The structure \texttt{HostConditions} looks like the analytic input of the
argument, including a pointwise bound for every bounded character. In the
archive it occurs only as a hypothesis: of
\texttt{EAtH.explicit\_formula\_error\_of\_host\_conditions} and of one
lemma in \arch{Development/AutomorphicArithmetic/CoprimeRanks/OfPInjective.lean}.
No term of this type is built on the path through \texttt{supplyChi}. That
path builds the explicit-formula condition \texttt{EAtH} directly
(Theorem~\ref{host:supply-of-adm-of}), with the pointwise and log-derivative
clauses supplied by Lemma~\ref{host:bounded-tw} and the Chebyshev clause
supplied inside the condition \texttt{HostGood}.
\end{remark}

% ------------------------------------------------------------------
\section{The host attached to a character}

\begin{definition}[kernel condition]\label{host:kerchi}
\lean{UpperTheory.Top.KerChi}
\uses{host:host-raw}
% uses-lean: Formalism.HC.PChar
\uses{an:pchar} % resolved
Let $\psi$ be a primitive character of conductor $N_{\psi}$ and $H$ a raw host
for $\phi$. Then $H$ satisfies the \emph{kernel condition for $\psi$} if
$N_{\psi}\mid Q_S$ and, for every $a$ prime to $N_H$, $\mathrm{art}(a)=1$ if
and only if $a\equiv\pm1\pmod{2^{e+2}}$, $\psi(a)=1$,
$\left(\frac aq\right)=1$ and $\left(\frac ad\right)=1$.
\end{definition}

\begin{definition}[field of $\psi$]\label{host:ischifield}
\lean{UpperTheory.Top.IsChiField}
% uses-lean: Formalism.HC.PChar
\uses{an:pchar} % resolved
A number field $E$ is a \emph{field of $\psi$ at $(e,q,d)$} if there are
$N\ge1$ with $2^{e+2},N_{\psi},q,d\mid N$, a $\Q$-embedding
$j\colon E\to\Q(\zeta_N)$ and a primitive $N$-th root of unity $\zeta$ such
that, for every automorphism $\tau$ with $\tau(\zeta)=\zeta^a$, $\tau$ fixes
$j(E)$ pointwise if and only if $a\equiv\pm1\pmod{2^{e+2}}$, $\psi(a)=1$,
$\left(\frac aq\right)=1$ and $\left(\frac ad\right)=1$.
\end{definition}

\begin{lemma}\label{host:ischifield-of-kerchi}
\lean{UpperTheory.Top.isChiField}
\uses{host:kerchi, host:ischifield}
If a raw host $H$ for $\phi$ satisfies the kernel condition for $\psi$, then
$E$ is a field of $\psi$ at $(e,q,d)$.
\end{lemma}
\begin{proof}
Take $N=N_H$, the embedding $j$ and the root $\zeta$ of the host. The
divisibilities follow from $2^{e+2}\mid Q_S$ and $N_{\psi}\mid Q_S$. If
$\tau(\zeta)=\zeta^a$ then $a$ is prime to $N_H$, and $\tau$ fixes $j(E)$
if and only if $\mathrm{art}(a)=1$, by the compatibility of $\mathrm{art}$
with $j$; the kernel condition finishes the proof.
\end{proof}

\begin{definition}[twist form]\label{host:twist-form}
\lean{UpperTheory.Top.TwistForm, UpperTheory.Top.tw}
\uses{host:chars, host:dir}
% uses-lean: Formalism.HC.PChar
\uses{an:pchar} % resolved
For $\xi\colon\N\to\CC$ put $\mathrm{tw}(\xi,q,d,i,j)(a)=\xi(a)\chi_q(a)^i
\chi_d(a)^j$, $i,j\in\{0,1\}$. A raw host $H$ has \emph{twist form for
$\psi$} if every character $\theta$ trivial on the first factor admits
$\xi\in\langle\psi\rangle$ and $i,j\in\{0,1\}$ with $\theta(\iota)=(-1)^j$
and $\theta(\mathrm{art}(a))=\mathrm{tw}(\xi,q,d,i,j)(a)$ for every $a$
prime to $N_H$.
\end{definition}

\begin{lemma}[the host of a cyclic group of characters]\label{host:any-grp}
\lean{UpperTheory.Top.ordAny, UpperTheory.Top.kerChiAny, UpperTheory.Top.twistAny, Formalism.HC.AnyGrp.host, Formalism.HC.AnyGrp.o_eq, Formalism.HC.AnyGrp.keptInXp, Formalism.HC.AnyGrp.exceptional_character_condition}
\uses{host:host-raw, host:kerchi, host:twist-form, host:kept}
% uses-lean: Formalism.HC.Level.ofSub
\uses{an:level-cyc} % resolved
Let $M\ge1$ and $\Xi$ a subgroup of the Dirichlet characters modulo $M$.
For every $e\ge2$ and every pair $q,d$ with \texttt{HostHyp}, $\gcd(M,qd)=1$
and $\left(\frac{-d}{l}\right)=1$ for all odd primes $l\mid M$, the archive
constructs a frame and a raw host $\mathrm{AnyGrp.host}(\Xi,e)$ with these
$q,d$ such that
\begin{itemize}
\item $o=|\Xi|$;
\item there are $1\le c\le|\Xi|$ and $e_0$ with $\mathrm{ord}=c\,2^e$ for
  all $e\ge e_0$;
\item it has kept characters in $X_+$ and keeps the exceptional character,
  for the kept set $\{\mathrm{prim}(\chi):\chi\in\Xi\}$ and empty
  exceptional slot.
\end{itemize}
For $\Xi=\langle\psi\rangle$ it moreover satisfies the kernel condition
for $\psi$ and has twist form for $\psi$.
\end{lemma}
\begin{proof}
The construction and the verification live in
\arch{Formalism/SharpRestriction/RamificationFrobenius/TotientGaloisClasses/InertiaGaloisSymmetry/AnyLevel.lean},
\arch{AnyFrame.lean} there, and
\arch{Development/AutomorphicArithmetic/CharacterBoxes/\{OrdAny,KerChiAny,TwistAny\}.lean};
they are not expanded here. The field is cut out by the subgroup of the
kernel condition.
\end{proof}

\begin{definition}[field rows]\label{host:field-rows}
\lean{HostParam.Seam.FieldRows}
\uses{host:host-raw}
A raw host $H$ has \emph{field rows with bound $b$} if: the first sign of
$\mathrm{art}(u)$ is $\left(\frac uq\right)$; there is $\sqrt{-q}\in E$ on
which $\sigma$ acts by the first sign; for every prime $l\mid Q_S$ some
$u\equiv1\pmod{N_H/l}$ has $\mathrm{art}(u)\ne1$; $Q_S\le\mathrm{ord}^2$;
$o^{64}\le2^e$; $q\ne3$; $\left(\frac{-d}{p}\right)=1$ for every odd prime
$p\le b$; and $4l\le b$ for every odd prime $l\mid Q_S$.
\end{definition}

\begin{theorem}[a host for every character]\label{host:field-chi}
\lean{UpperTheory.Top.fieldChi}
\uses{host:any-grp, host:host-ax, host:field-rows, host:kept, host:kerchi, host:twist-form}
% uses-lean: Formalism.HC.Level.cyc, Formalism.HC.Frame.two_pow_le_ord
\uses{an:level-cyc} % resolved
Let $\psi$ be a primitive character of conductor $N_{\psi}$. There are
integers $c\ge1$ and $e_0$ such that the following holds for all $e\ge e_0$,
all primes $q,d$ with \texttt{HostHyp}, $q\ne3$ and $\gcd(N_{\psi},qd)=1$, and
all $b\ge4N_{\psi}$ with $\left(\frac{-d}{p}\right)=1$ for every odd prime
$p\le b$. There are a frame $\phi$ and a host with axioms $H$ for $\phi$
with $e(\phi)=e$, $\mathrm{ord}=c\,2^e$, auxiliary primes $q,d$,
$o(\phi)=|\langle\psi\rangle|$, field rows with bound $b$,
$Q_S\mid\operatorname{lcm}(2^{e+2},N_{\psi})$, $Q_S\le4N_{\psi}\,\mathrm{ord}$,
kept characters in $X_+$ and the exceptional character kept for the kept
set of powers of $\psi$ (with empty exceptional slot), the kernel condition
for $\psi$ and twist form for $\psi$.
\end{theorem}
\begin{proof}
\uses{host:any-grp, host:host-ax}
Apply Lemma~\ref{host:any-grp} to $\Xi=\langle\psi\rangle$, with $e_0$
the maximum of the threshold for $\mathrm{ord}=c2^e$, of $2$, of
$4N_{\psi}$ and of $64|\langle\psi\rangle|$. The splitting hypothesis for
primes dividing $N_{\psi}$ follows from the one up to $b\ge4N_{\psi}$. Upgrade
the raw host with \texttt{HostAx.ofRaw}. For the field rows:
$o^{64}=|\langle\psi\rangle|^{64}\le2^{64|\langle\psi\rangle|}\le2^e$;
every odd prime $l\mid Q_S$ divides $N_{\psi}$, so $4l\le4N_{\psi}\le b$; and
$Q_S\le4N_{\psi}2^e\le4N_{\psi}\,\mathrm{ord}$ because $2^e\le\mathrm{ord}$.
\end{proof}

\begin{remark}[the exceptional slot is empty]\label{host:rem-exc}
For the level of powers of $\psi$ the exceptional slot is empty, so
``keeps the exceptional character'' says only that $\omega_q$, whose
primitive character is the quadratic character attached to $q$, is not
the primitive character of a power of $\psi$.
\end{remark}

% ------------------------------------------------------------------
\section{Window pairs and good pairs}

\begin{definition}[window pairs]\label{host:win-pair}
\lean{UpperTheory.Top.WinPair, UpperTheory.Top.PageN}
Let $\eps,Y>0$. A pair of integers $(q,d)$ is a \emph{window pair} for
$(\eps,Y)$ if $q$ is a prime with $q\equiv3\pmod8$ and
$Y^{\eps/5}/2<\log q\le Y^{\eps/5}$; $d$ is a prime with
$d\equiv7\pmod{16}$ and $Y^{\eps/5}\log^2Y/2<\log d\le Y^{\eps/5}\log^2Y$;
$q\mid d+1$; and $p\mid d+1$ for every odd prime $p\le Y^{\eps/5}$. An
integer $e$ is \emph{adapted} to $(\eps,x)$ if
$x^{2+8\eps}\le2^e\le2x^{2+8\eps}$.
\end{definition}

\begin{lemma}\label{host:exists-pagen}
\lean{UpperTheory.Top.exists_pageN}
\uses{host:win-pair}
For every $\eps>0$ there is a function $x\mapsto e(x)$ with $e(x)$ adapted
to $(\eps,x)$ for all $x\ge1$.
\end{lemma}
\begin{proof}
For $x\ge1$, $x^{2+8\eps}\ge1$, so there is $n$ with
$2^n\le x^{2+8\eps}<2^{n+1}$; take $e(x)=n+1$.
\end{proof}

\begin{lemma}[window pairs give host primes]\label{host:pair-hyp}
\lean{UpperTheory.Top.pairHyp}
\uses{host:win-pair, host:host-raw}
Let $(q,d)$ be a window pair for $(\eps,Y)$. Then $(q,d)$ satisfies
\texttt{HostHyp}; $\left(\frac{-d}{p}\right)=1$ for every odd prime
$p\le Y^{\eps/5}$; and $\left(\frac{-q}{a}\right)=\left(\frac aq\right)$ and
$\left(\frac{-d}{a}\right)=\left(\frac ad\right)$ for every odd $a>0$.
\end{lemma}
\begin{proof}
If $p\mid d+1$ then $-d\equiv1\pmod p$, so $\left(\frac{-d}{p}\right)=1$;
this gives $\left(\frac{-d}{q}\right)=1$ and the splitting of small primes.
The last identities are quadratic reciprocity for the Jacobi symbol, using
$q\equiv d\equiv3\pmod4$.
\end{proof}

\begin{definition}[good pairs]\label{host:pair-good}
\lean{UpperTheory.Top.PairGood, UpperTheory.Top.ZeroBox, UpperTheory.Top.height, UpperTheory.Top.PairEx}
\uses{host:win-pair, host:twist-form}
% uses-lean: Formalism.HC.An.Induces, Formalism.HC.An.gChi, Formalism.HC.An.logDerivativeAtOne, Formalism.HC.PChar.IsReal
\uses{an:gchi, an:pchar} % resolved
Let $M\ge1$, $\eps>0$ and $Y\le X$. A Dirichlet character $\chi$ has a
\emph{zero-free box of height $Q$} if $L(s,\chi)\ne0$ whenever
$\Re s>1-1/400$ and $|\Im s|\le Q$. Put
$H_j=8M\exp(Y^{\eps/5})\exp(Y^{\eps/5}\log^2Y)^j$ for $j\in\{0,1\}$. A
pair $(q,d)$ is \emph{good} for $(M,\eps,Y,X)$ if it is a window pair for
$(\eps,Y)$, $\gcd(M,qd)=1$, and
\begin{enumerate}
\item for every Dirichlet character $\xi$ modulo $M$, every $(i,j)\ne(0,0)$
  and every primitive $\chi$ agreeing with $\mathrm{tw}(\xi,q,d,i,j)$ on
  integers prime to $2Mqd$, and every $y\in[Y,X]$:
  $|g_{\chi}(y)|\le y^{1/2+\eps/2}$ and $|L'(1,\chi)/L(1,\chi)|\le4\log y$;
\item for every real primitive $\varrho$ with conductor dividing $8M$,
  every $(i,j)\ne(0,0)$ and every primitive $\chi$ agreeing with
  $\mathrm{tw}(\varrho,q,d,i,j)$ on integers prime to $2Mqd$, $\chi$ has a
  zero-free box of height $H_j$.
\end{enumerate}
The proposition \texttt{PairEx} states: for every $M>0$, every
$\eps\in(0,1/52)$ and every $A\ge1$ there is $Y_1$ such that for every
$Y\ge Y_1$ some pair is good for $(M,\eps,Y,Y^A)$.
\end{definition}

\begin{lemma}[bounded characters are proper twists]\label{host:bounded-tw-key}
\lean{UpperTheory.Top.boundedTw_key}
\uses{host:twist-form, host:kept, host:prim-unique}
% uses-lean: Formalism.HC.Level.cyc, Formalism.HC.Level.mem_ofSub, Formalism.HC.An.induces_of
\uses{an:level-cyc} % resolved
Let $\psi$ be primitive of conductor $N_{\psi}$ and $H$ a raw host with twist
form for $\psi$, $N_{\psi}\mid Q_S$ and every prime divisor of $Q_S$ dividing
$2N_{\psi}$. Let $S$ be the set of primitive characters of the powers of
$\psi$. If $\theta$ is bounded for $S$ and $\chi$ is the primitive
character of $\theta$, then there are a character $\xi$ modulo $N_{\psi}$ and
$(i,j)\ne(0,0)$ such that $\chi$ agrees with $\mathrm{tw}(\xi,q,d,i,j)$ on
integers prime to $2N_{\psi} qd$.
\end{lemma}
\begin{proof}
A bounded $\theta$ is trivial on the first factor, so twist form gives
$\xi\in\langle\psi\rangle$ and $i,j$. If $(i,j)=(0,0)$, then either
$\theta\in X_-$, contradicting $\theta(\iota)=(-1)^0=1$, or $\theta\in X_+$
is not kept; but then $\chi$ agrees with $\xi$ on integers prime to $N_H$,
so $\chi=\mathrm{prim}(\xi)\in S$ and $\theta$ is kept, a contradiction.
For the agreement: an integer $a$ prime to $2N_{\psi} qd$ is prime to $Q_S$,
since every prime divisor of $Q_S$ divides $2N_{\psi}$, hence prime to $N_H$,
and there $\chi(a)=\theta(\mathrm{art}(a))=\mathrm{tw}(\xi,q,d,i,j)(a)$.
\end{proof}

\begin{lemma}[bounds for the bounded characters]\label{host:bounded-tw}
\lean{UpperTheory.Top.boundedTw}
\uses{host:bounded-tw-key, host:pair-good, host:kept}
Let $\psi$, $H$ and $S$ be as in Lemma~\ref{host:bounded-tw-key}, let
$\eps>0$, $x\ge1$, $Y\le x$, and assume that the auxiliary primes $(q,d)$
of $H$ are good for $(N_{\psi},\eps,Y,x)$. Then $H$ has pointwise bounds at
$x$ with exponent $\eps$ for $S$, and $|L'(1,\chi)/L(1,\chi)|\le4\log x$
for the primitive character $\chi$ of every bounded $\theta$.
\end{lemma}
\begin{proof}
\uses{host:bounded-tw-key}
By Lemma~\ref{host:bounded-tw-key}, $\chi$ is a proper twist to which
clause (1) of Definition~\ref{host:pair-good} applies at $y=x\in[Y,x]$.
This gives $|g_{\chi}(x)|\le x^{1/2+\eps/2}\le x^{1/2+\eps}$, as $x\ge1$, and
the log-derivative bound.
\end{proof}

% ------------------------------------------------------------------
\section{Averaging over the auxiliary primes}

\begin{definition}[head sums and prime families]\label{host:head-sum}
\lean{QH.S, QH.J, QH.MeanSq, QH.admitted, Formalism.HermitianLattices.ZeroFreeRegionGL2.cN, Formalism.HermitianLattices.ZeroFreeRegionGL2.PPset, Formalism.HermitianLattices.ZeroFreeRegionGL2.oddM}
Let $\Lambda$ be the von Mangoldt function.
\begin{enumerate}
\item For $Z,y>0$, an integer $q\ge1$ and $\kappa\colon\N\to\CC$, the
  \emph{head sum} is
  $S(Z,y,q,\kappa)=\sum_{1\le n\le Z}\Lambda(n)\min(1,y/n)\,\kappa(n)
  \left(\frac nq\right)$.
\item For integers $m_0,a_0$ and $U>0$, $P(m_0,a_0,U)$ is the set of primes
  $q$ with $U^{1/2}<q\le U$ and $q\equiv a_0\pmod{m_0}$.
\item For $y>0$ and an integer $q$, $\mathrm{oddM}(y,q)=q\prod_{3\le p\le
  y}p$, the product over primes.
\item $(m_0,a_0)$ is \emph{admitted} for $(\eps,Y,U)$ if $m_0>0$, $8\mid
  m_0$, $\gcd(a_0,m_0)=1$, $\frac14(\log Y)^2\log m_0\le\log U$ and
  $Y^{\eps/6}\le\log U\le Y^{\eps/5}(\log Y)^2$.
\item For $C,Z\in\R$, a finite set $P$ and $I\subset\R$, the
  \emph{mean-square bound} $\mathrm{MeanSq}(C,Z,P,I)$ holds if
  $\sum_{q\in P}|S(Z,y,q,\kappa)|^2\le C\,y\log y\,|P|$ for all $y\in I$
  and all $\kappa$ with $|\kappa(n)|\le1$.
\end{enumerate}
\end{definition}

\begin{definition}[boxes]\label{host:box}
\lean{QH.Ind, QH.GoodC, QH.Box, QH.BoxAll, Formalism.HermitianLattices.ZeroFreeRegionGL2.ZeroFreeBox, Formalism.HermitianLattices.ZeroFreeRegionGL2.Gchi, Formalism.HermitianLattices.ZeroFreeRegionGL2.logDerivLAtOne}
\uses{host:pair-good}
Let $\chi$ be a Dirichlet character modulo $N$, and let
$G_{\chi}(y)=\sum_{n\le y}\Lambda(n)\chi(n)(y/n-1)+y\,L'(1,\chi)/L(1,\chi)$.
\begin{enumerate}
\item $\chi$ is \emph{good at height $Q$} if $\chi\ne1$, $N\le Q$ and
  $\chi$ has a zero-free box of height $Q$.
\item $\chi$ \emph{has a box} $\mathrm{Box}(Q,C,\tau,Y,X)$ if it is good at
  height $Q$, $|L'(1,\chi)/L(1,\chi)|\le C\log Y$ and
  $|G_{\chi}(y)|\le y^{1/2+\tau}$ for all $y\in[Y,X]$.
\item For $f\colon\N\to\CC$ and $M$, $\mathrm{BoxAll}(Q,C,\tau,Y,X,f,M)$
  holds if every primitive $\chi$ agreeing with $f$ on integers prime to
  $M$ has a box $\mathrm{Box}(Q,C,\tau,Y,X)$.
\end{enumerate}
\end{definition}

\begin{definition}[good auxiliary primes]\label{host:good-qd}
\lean{QH.GoodQ, QH.GoodD, QH.Rows, QH.primePairSetFirst, QH.primeSetTwoBounds}
\uses{host:box, host:head-sum}
Let $N_0\ge1$, $C,\tau,Y,X,z,z'\in\R$ and $J_q(n)=\left(\frac nq\right)$.
\begin{enumerate}
\item $q$ is \emph{good} if $\mathrm{BoxAll}(N_0e^z,C,\tau,Y,X,\xi J_q,
  N_0q)$ for every character $\xi$ modulo $N_0$.
\item $d$ is \emph{good for $q$} if $\mathrm{BoxAll}(N_0e^{z+z'},C,\tau,Y,X,
  \xi J_q^aJ_d,N_0qd)$ for every $\xi$ modulo $N_0$ and $a\in\{0,1\}$.
\item The first family is $P(8,3,e^z)$. For $q$, the second family is the
  set of primes $d\le e^{z'}$ with $d>e^{z'/2}$, $d\equiv7\pmod{16}$ and
  $\mathrm{oddM}(z,q)\mid d+1$.
\item $(q,d)$ are \emph{rows} for $(z,z')$ if $q$ is a prime, $q\equiv3
  \pmod8$, $z/2<\log q\le z$; $d$ is a prime, $d\equiv7\pmod{16}$,
  $z'/2<\log d\le z'$; $q<d$; and $p\mid d+1$ for every odd prime $p$ with
  $p\le z$ or $p=q$.
\end{enumerate}
\end{definition}

\begin{lemma}[mean square over prime moduli]\label{host:mean-sq}
\lean{QH.meanSqK}
\uses{host:head-sum}
There is $C>0$ such that for all $\eps\in(0,1]$ and $K\ge1$ there is $Y_0$
with the following property. Let $Y\ge Y_0$, $U>0$, let $(m_0,a_0)$ be
admitted for $(\eps,Y,U)$, and assume every prime $p\mid m_0$ satisfies
$p\le Y^{\eps/5}$ or $p>Y^{400K}$. Then
$\mathrm{MeanSq}(C,Y^{400K},P(m_0,a_0,U),[Y,Y^K])$ holds.
\end{lemma}
\begin{proof}
The proof is in \arch{Development/AutomorphicArithmetic/MeanSquareBounds/meanSqK.lean}
and its imports (\texttt{QH.twistK}, a prime number theorem in progressions
with relative error); it was not read here.
\end{proof}

\begin{lemma}[few bad moduli for the head sum]\label{host:bad-share}
\lean{QH.bad_share}
\uses{host:head-sum}
Let $C\ge0$, $\tau>0$ and $\eta>0$. There is $Y_0$ such that for all
$Y\ge Y_0$, all $X,Z$, all finite sets $B\subseteq P$ of natural numbers and
all $\kappa$ with $|\kappa(n)|\le1$: if $\mathrm{MeanSq}(C,Z,P,[Y,X])$ holds
and every $q\in B$ has some $y\in[Y,X]$ with
$|S(Z,y,q,\kappa)|>y^{1/2+\tau}/2$, then $|B|\le\eta|P|$.
\end{lemma}
\begin{proof}
In \arch{Development/AutomorphicArithmetic/MeanSquareBounds/BadShare.lean};
not read here.
\end{proof}

\begin{lemma}[few moduli without a zero-free box]\label{host:count-bad}
\lean{QH.count_bad}
\uses{host:pair-good}
% uses-lean: Formalism.HermitianLattices.Ext.PrimitiveCharacterCasts.ZD.zero_density_box, QH.count_bad_core
There is $C$ such that for all $Q\ge3$, all $F$ and every finite set $B$ of
natural numbers the following holds. Assume each $q\in B$ is a prime with
$q>F$ for which there are an integer $f\le F$ with $fq\le Q$ and a primitive
character modulo $fq$ without a zero-free box of height $Q$. Then
$|B|\le CQ^{1/4}$.
\end{lemma}
\begin{proof}
From \texttt{count\_bad\_core}, which rests on the zero-density statement
\texttt{ZD.zero\_density\_box}; not read here.
\end{proof}

\begin{lemma}[small head sum gives a box]\label{host:box-of-s}
\lean{QH.box_of_S}
\uses{host:box, host:head-sum}
There is $C>0$ such that for all $\eps\in(0,1/4]$, $\tau>0$, $K\ge1$ there
is $Y_0$ with the following property. Let $Y\ge Y_0$, let $\chi$ be a
character modulo $N$, $q$ an integer and $\kappa\colon\N\to\CC$ with
$\chi(n)=\kappa(n)\left(\frac nq\right)$ for all $n$, and let $Q$ satisfy
$Y^{\eps/6}\le\log Q\le Y^{\eps/4}$. If $\chi$ is good at height $Q$ and
$|S(Y^{400K},y,q,\kappa)|\le y^{1/2+\tau}/2$ for all $y\in[Y,Y^K]$, then
$\chi$ has a box $\mathrm{Box}(Q,C,\tau,Y,Y^K)$.
\end{lemma}
\begin{proof}
% uses-lean: QH.kappaC, QH.tailK
Uses \texttt{QH.kappaC} (a zero-free box bounds $L'(1,\chi)/L(1,\chi)$ by
$C\log Y$) and a tail estimate showing that the part of $G_{\chi}(y)$ beyond
$Y^{400K}$ is negligible under the zero-free box; proof in
\arch{Development/AutomorphicArithmetic/MeanSquareBounds/BoxOfS.lean}, not
read in detail.
\end{proof}

\begin{lemma}[the constant four]\label{host:kappa4}
\lean{QH.kappa4}
\uses{host:box}
For all $\eps\in(0,1/52)$ and $A>0$ there is $Y_0$ such that for all
$Y\ge Y_0$, all $Q$ with $Y^{\eps/6}\le\log Q\le AY^{\eps/5}\log^2Y$ and
every character $\chi$ good at height $Q$,
$|L'(1,\chi)/L(1,\chi)|\le4\log Y$.
\end{lemma}
\begin{proof}
From \texttt{kappa4\_loglog} in
\arch{Development/AutomorphicArithmetic/MeanSquareBounds/Kappa4Loglog.lean};
not read here.
\end{proof}

\begin{lemma}[one stage of the averaging]\label{host:stage}
\lean{QH.stage}
\uses{host:mean-sq, host:bad-share, host:count-bad, host:box-of-s, host:head-sum, host:box}
% uses-lean: Formalism.HermitianLattices.Ext.LogarithmicIntegralBounds.count_a, Formalism.HermitianLattices.Ext.PrimesInProgressions.pnt_progressions_relative_weak_holds, Formalism.ZeroFreeDirichlet.excBound
There is $C>0$ such that for all $\eps\in(0,1/4]$, $\tau>0$, $K\ge1$,
$\eta>0$ and $R\in\N$ there is $Y_0$ with the following property. Let
$Y\ge Y_0$, $(m_0,a_0)$ admitted for $(\eps,Y,U)$ with every prime
$p\mid m_0$ satisfying $p\le Y^{\eps/5}$ or $p>Y^{400K}$, let $1\le F\le
U^{1/2}$, and let $\Xi$ be a set of at most $R$ primitive characters of
conductor at most $F$. Then $P(m_0,a_0,U)$ is nonempty, and every subset
$B\subseteq P(m_0,a_0,U)$ such that each $q\in B$ has some $\xi\in\Xi$ and
$M\ne0$ with $\mathrm{BoxAll}(FU,C,\tau,Y,Y^K,\xi J_q,M)$ false satisfies
$|B|\le\eta\,|P(m_0,a_0,U)|$.
\end{lemma}
\begin{proof}
\uses{host:mean-sq, host:bad-share, host:count-bad, host:box-of-s}
By a prime number theorem in progressions (\texttt{count\_a}),
$|P(m_0,a_0,U)|\ge U^{9/10}$. For $q\in B$ and the offending $\xi$, the
primitive character $\chi$ inducing $\xi J_q$ has conductor
$\mathrm{cond}(\xi)\,q\le FU$. Split $B$ into $B_1$, the $q$ for which
some primitive character of conductor $fq\le FU$ with $f\le F$ has no
zero-free box of height $FU$, and for each $\xi\in\Xi$ the set $B_2(\xi)$
of $q$ for which $|S(Y^{400K},y,q,\xi)|>y^{1/2+\tau}/2$ for some
$y\in[Y,Y^K]$. If $q$ is in neither, Lemma~\ref{host:box-of-s} gives
$\chi$ a box, and by uniqueness of primitive characters
$\mathrm{BoxAll}$ holds, so $B\subseteq B_1\cup\bigcup_{\xi} B_2(\xi)$.
Lemma~\ref{host:count-bad} gives $|B_1|\le C_b(FU)^{1/4}\le
C_bU^{3/8}\le\frac\eta2U^{9/10}$ for $U$ large. Lemma~\ref{host:mean-sq}
and Lemma~\ref{host:bad-share} (with $\eta/(2(R+1))$) give
$|B_2(\xi)|\le\frac{\eta}{2(R+1)}|P|$, and summing over at most $R$
characters gives $|\bigcup B_2(\xi)|\le\frac\eta2|P|$.
\end{proof}

\begin{lemma}[the two stages]\label{host:stages}
\lean{QH.stage1, QH.stage2}
\uses{host:stage, host:good-qd}
There is $C>0$ such that for all $\eps\in(0,1/52)$, $\tau>0$, $K\ge1$,
$\eta>0$ and $N_0>0$ there is $Y_0$ such that for $Y\ge Y_0$ and
$z=Y^{\eps/5}$:
\begin{enumerate}
\item the first family $P(8,3,e^z)$ is nonempty and each subset of it
  consisting of primes $q$ that are not good (for $(N_0,C,\tau,Y,Y^K,z)$)
  has at most $\eta$ times its size;
\item for every $z'$ with $z\log^2Y\le z'\le K^2z\log^2Y$ and every $q$ in
  the first family, the second family for $q$ is nonempty and each subset
  of it consisting of $d$ not good for $q$ has at most $\eta$ times its
  size.
\end{enumerate}
\end{lemma}
\begin{proof}
\uses{host:stage}
Both are instances of Lemma~\ref{host:stage}: the first with
$(m_0,a_0)=(8,3)$, $U=e^z$, the second with $m_0=16\,\mathrm{oddM}(z,q)$,
$U=e^{z'}$, exponent $6\eps/5$ in place of $\eps$, a residue class $a_0$
realising the second family, and the characters $\xi J_q^a$. The prime
divisors of $16\,\mathrm{oddM}(z,q)$ are $2$, odd primes at most $z$, and
$q$, which exceeds $Y^{400K}$ for $Y$ large.
The proofs are in \arch{.../MeanSquareBounds/Stage1Succ.lean} and
\arch{.../MeanSquareBounds/CharOf.lean}; only their opening lines were read.
\end{proof}

\begin{lemma}[most pairs are good]\label{host:pair-share}
\lean{QH.pair_share}
\uses{host:stages}
There is $C>0$ such that for all $\eps\in(0,1/52)$, $\tau>0$, $K\ge1$,
$\eta>0$ and $N_0>0$ there is $Y_0$ such that for all $Y\ge Y_0$, $X\le
Y^K$ and $z'$ with $z\log^2Y\le z'\le K^2z\log^2Y$, $z=Y^{\eps/5}$: there is
a subset $G_1$ of the first family, of size at least $(1-\eta)$ times its
positive size, all of whose elements $q$ are good for $(N_0,C,\tau,Y,X,z)$
and admit a subset $G_2$ of the second family for $q$, of size at least
$(1-\eta)$ times its positive size, all of whose elements are good for $q$.
\end{lemma}
\begin{proof}
\uses{host:stages}
Take $G_1$ (resp.\ $G_2$) to be the good elements; the complement is a bad
set, bounded by Lemma~\ref{host:stages}. Monotonicity of $\mathrm{BoxAll}$
in $X$ passes from $Y^K$ to $X\le Y^K$.
\end{proof}

\begin{lemma}[existence of a good pair of rows]\label{host:pair-exists}
\lean{QH.pair_exists}
\uses{host:pair-share, host:good-qd}
% uses-lean: QH.rows_of_mem
There is $C>0$ such that for all $\eps\in(0,1/52)$, $\tau>0$, $K\ge1$ and
$N_0>0$ there is $Y_0$ such that for all $Y\ge Y_0$, $X\le Y^K$ and $z'$
with $z\log^2Y\le z'\le K^2z\log^2Y$, $z=Y^{\eps/5}$, there are rows $(q,d)$
for $(z,z')$ with $q>N_0$, $q$ good and $d$ good for $q$.
\end{lemma}
\begin{proof}
\uses{host:pair-share}
Apply Lemma~\ref{host:pair-share} with $\eta=1/2$. Both good sets are
nonempty because they have at least half of a positive size; pick $q\in
G_1$, then $d\in G_2$. Membership in the two families gives rows. For $Y$
large, $2N_0\le Y^{\eps/5}<2\log q$, so $q>N_0$.
\end{proof}

\begin{theorem}[good pairs exist]\label{host:pair-ex}
\lean{QH.pairEx}
\uses{host:pair-exists, host:kappa4, host:pair-good, host:good-qd}
The proposition \texttt{PairEx} of Definition~\ref{host:pair-good} holds:
for every $M>0$, $\eps\in(0,1/52)$ and $A\ge1$ there is $Y_1$ such that for
every $Y\ge Y_1$ some pair is good for $(M,\eps,Y,Y^A)$.
\end{theorem}
\begin{proof}
\uses{host:pair-exists, host:kappa4}
Apply Lemma~\ref{host:pair-exists} with $\tau=\eps/2$, $K=A$, $N_0=8M$
and $z'=Y^{\eps/5}\log^2Y$, and Lemma~\ref{host:kappa4} with $A=3$. Rows
give a window pair: $q\mid d+1$ and the small primes divide $d+1$. Since
$q,d>8M$ are primes, $\gcd(M,qd)=1$. For clause (1), lift $\xi$ to modulus
$8M$; goodness of $q$ (for $j=0$) or of $d$ for $q$ (for $j=1$) gives every
primitive $\chi$ inducing the twist a box of height
$8Me^{z}e^{z'j}=H_j$, hence $|g_{\chi}(y)|\le y^{1/2+\eps/2}$ on $[Y,Y^A]$.
The height satisfies $Y^{\eps/6}\le\log H_j\le3Y^{\eps/5}\log^2Y$, so
Lemma~\ref{host:kappa4} gives $|L'(1,\chi)/L(1,\chi)|\le4\log Y\le4\log y$.
Clause (2) is the zero-free-box part of the same boxes, applied to the
real character $\varrho$ lifted to modulus $8M$.
\end{proof}

\begin{remark}[where the constant four comes from]\label{host:rem-four}
The constant $4$ in $|L'(1,\chi)/L(1,\chi)|\le4\log y$ is not produced by
Lemma~\ref{host:box-of-s}, whose constant $C$ is unspecified. It is
inserted in the proof of Theorem~\ref{host:pair-ex} by
Lemma~\ref{host:kappa4} with $A=3$.
\end{remark}

\begin{remark}[character of the averaging]\label{host:rem-averaging}
Lemma~\ref{host:bounded-tw} and Theorem~\ref{host:pair-ex} need no
hypothesis of Riemann type. The scales $y\in[Y,Y^A]$ satisfy $\log q\asymp
Y^{\eps/5}$, so $y$ is polylogarithmic in the conductor, and the head sums
have length $Y^{400A}$, far below $q$. The zero-free boxes reach a height
that is a power of the conductor, but only in the region
$\Re s>1-1/400$. Which classical theorems the unread proofs of
Lemmas~\ref{host:mean-sq}--\ref{host:kappa4} use was not determined here;
\texttt{QH.stage} names the prime number theorem in progressions with
relative error and the Siegel-type bound \texttt{excBound}.
\end{remark}

% ------------------------------------------------------------------
\section{The supply theorem}

\begin{definition}[scales and windows]\label{host:scale}
\lean{Formalism.BasicNotions.Scale, GlobalTheory.EX.Scales, GlobalTheory.EX.Scales.Mem, GlobalTheory.EX.PairWin}
A \emph{scale} $s$ is $(x,\eps,J)$ with $x>1$, $0<\eps<1/52$ and $J\ge1$;
put $b=\exp(x^{\eps})$. A \emph{system of windows} of exponent $\lambda$ is
$J\ge1$ and a set $W\subset\R$ such that for every $y$ there is $Y\ge y$
with $[Y,Y^{\lambda}]\subseteq W$; a scale is a \emph{member} if its $J$ is
the system's $J$ and $x\in W$. A real $Y$ is a \emph{pair window} for $s$ if
$x^{1/J}\le Y\le x$.
\end{definition}

\begin{definition}[rank values]\label{host:rank-val}
\lean{GlobalTheory.EX.RankVal}
\uses{host:ranks, host:scale}
A rank datum $R$ has \emph{values for the scale $s$} if
$x^{2+8\eps}\le n\le2x^{2+8\eps}$, $\sqrt{x\delta_1}\le m\le2\sqrt{x\delta_1}$,
$x^{9/2+10\eps}\le k\le2x^{9/2+10\eps}$, and $r+t$ equals an explicit
expression in $f,r$ and $s$ (here $\delta_1$ is a quantity attached to $s$).
\end{definition}

\begin{definition}[good host]\label{host:host-good}
\lean{GlobalTheory.EX.HostGood, GlobalTheory.EX.HostVal, GlobalTheory.EX.ThinOK, GlobalTheory.EX.BoxW, GlobalTheory.EX.dlogH}
\uses{host:host-full, host:scale, host:kept}
% uses-lean: Formalism.MainNotions.Inhab.QDq, GlobalTheory.EX.BadOK, GlobalTheory.EX.HostArith
For a scale $s$, constants $C_8,C_5,c_6,C_b$, a host $H$ for $R$ and an
integer $q_0$, $\mathrm{HostGood}(s,C_8,C_5,c_6,C_b,H,q_0)$ is the
conjunction of: size conditions (a pair window $Y$ for which $q,d$ lie in
the ranges $(\sqrt{X_1},X_1]$, $(\sqrt{X_2},X_2]$ with
$X_1=\exp(Y^{\eps/5})$, $X_2=\exp(Y^{\eps/5}\log^2Y)$, small primes split
in $K$, $\log N_H\le2Y^{\eps/5}\log^2Y$; sizes of $TT$, $T_2$, $g_i$,
$p_{\delta}$; $J\le C_8$, $o\le C_8$, $Q_S\le C_8n$); arithmetic of the level
primes; bounds with constant $C_5$ for the three real characters attached
to $q$, $d$ and $qd$ at $x$; zero-free boxes for the real characters of $E$
whose conductor is divisible by $q$ or $d$; $|\Re\,L'(1,\chi_{-d})/
L(1,\chi_{-d})|\le x^{\eps}$; the Chebyshev bound $(x,\eps,c_6)$ for every
character of $E$ (\texttt{ThinOK}, equivalent to
\texttt{ChebyshevCharacterBound} by \texttt{thinScales\_iff\_hostBound});
a count of bad Frobenius classes with constant $C_b$; and congruence and
size conditions making $q_0$ a prime $\equiv1\pmod{N_H}$ with $q_0\le
N_H^{C_8}$, and fixing the level primes as $\{g_0,g_1,p_{\delta},p'_{\delta},
q_0\}$.
\end{definition}

\begin{definition}[host supply]\label{host:host-exw}
\lean{GlobalTheory.EX.HostExW, GlobalTheory.EX.HostReq, GlobalTheory.EX.HostEx}
\uses{host:host-good, host:rank-val, host:scale}
A \emph{requirement} $P$ is a predicate on $(C_5,s,R,H,q_0)$. For
$\lambda\in\R$, $\mathrm{HostExW}(\lambda,P)$ states: for every
$\eps\in(0,1/52)$ there are constants $C_{T2}$, $C_8\ge1$, $C_5>0$,
$c_6>0$, $C_b>0$, $x_0$ and a system of windows of exponent $\lambda$ such
that every member scale $s$ with $\eps(s)=\eps$ and $x\ge x_0$ admits a rank
datum $R$ with values for $s$, a host $H$ for $R$ with
$C_{T2}(H)=C_{T2}$ and an integer $q_0$ with
$\mathrm{HostGood}(s,C_8,C_5,c_6,C_b,H,q_0)$ and $P(C_5,s,R,H,q_0)$.
\texttt{HostEx} is the variant with some $\lambda>1$ and $P$ always true.
\end{definition}

\begin{definition}[the requirement of a character]\label{host:gchi}
\lean{UpperTheory.Top.GChi, UpperTheory.Top.FReq.req}
\uses{host:host-exw}
% uses-lean: Formalism.HC.EAtH, Formalism.HC.Level.cyc, UpperTheory.Top.VVle
\uses{an:level-cyc, sha:vwin} % resolved
For a primitive character $\psi$, the requirement $\mathrm{GChi}(\psi)$
holds at $(s,R,H)$ if $H$ satisfies the explicit-formula condition
\texttt{EAtH} for the level of powers of $\psi$, the set $TT$ of $H$, the
point $x$, the exponent $\eps$ and the constant $4$; and the field $E$ of
$H$ satisfies the Sha-length bound \texttt{VVle} with $t(R)$ twists and
budget $\sqrt b$.
\end{definition}

\begin{definition}[window datum]\label{host:win-datum}
\lean{UpperTheory.Top.WinDatum}
\uses{host:pair-good, host:ischifield}
% uses-lean: UpperTheory.Top.VVle
\uses{sha:vwin} % resolved
For $\psi$, reals $\eps,J,c,x,Y$ and integers $q,d,e$, a \emph{window
datum} consists of: $x^{1/J}\le Y\le x$; $(q,d)$ good for
$(N_{\psi},\eps,Y,x)$; $x^{2+8\eps}\le c2^e\le2x^{2+8\eps}$; and the bound
\texttt{VVle} with $\lfloor x^{400}\rfloor$ twists and budget
$\sqrt{\exp(x^{\eps})}$ for every field of $\psi$ at $(e,q,d)$.
\end{definition}

\begin{lemma}[window data exist on windows]\label{host:win-pairs}
\lean{UpperTheory.Top.winPairs}
\uses{host:win-datum, host:exists-pagen, host:pair-good}
% uses-lean: UpperTheory.Top.VWinAt, UpperTheory.Top.VVle.mono, UpperTheory.Top.PairGood.mono
\uses{sha:vwin} % resolved
Let $\psi$ be primitive, $\eps\in(0,1/52)$ and $c\ge1$. Assume
\texttt{VWinAt}$(\psi,401)$ and \texttt{PairEx}. Then there is $J\ge1$ such
that for every $y$ there is $Z\ge y$ with: every $x\in[Z,Z^{3/2}]$ admits
$Y,q,d,e$ forming a window datum for $(\psi,\eps,J,c,x)$.
\end{lemma}
\begin{proof}
\uses{host:exists-pagen}
Fix an adapted function $e(\cdot)$ (Lemma~\ref{host:exists-pagen}), and
let \texttt{VWinAt} give $J$ and $X_0$. Write $c=k^{2+8\eps}$ with $k\ge1$.
Given $y$, choose $M$ above $y$, $1$, the threshold $Y_1$ of \texttt{PairEx}
(for $N_{\psi}$, $\eps$, exponent $J$), $k^3$ and $k^{400}$, and $X\ge X_0$
with $X\ge M^J$. \texttt{PairEx} at $X^{1/J}$ gives a pair good for
$(N_{\psi},\eps,X^{1/J},X)$, in particular a window pair, so \texttt{VWinAt}
gives $W\in[X^{1/J},X^{1/2}]$ with the Sha bound for all $x'\in[W,W^2]$.
Put $Z=kW$. Since $k^3\le W$, $(kW)^{3/2}\le W^2$. For $x\in[Z,Z^{3/2}]$
write $x=kx'$ with $x'\in[W,W^2]$, and take $Y=X^{1/J}$ and $e=e(x')$.
Then $c2^{e}$ lies between $x^{2+8\eps}$ and $2x^{2+8\eps}$, the good pair
restricts from $X$ to $x$, and the Sha bound at $x'$ with $\lfloor
x'^{401}\rfloor\ge\lfloor x^{400}\rfloor$ twists (as $k^{400}\le x'$) and
budget $\sqrt{\exp(x'^{\eps})}\le\sqrt{\exp(x^{\eps})}$ gives the required
one by monotonicity.
\end{proof}

\begin{lemma}\label{host:rank-le}
\lean{UpperTheory.Top.eventually_rank_le_rpow_of_lt_one_div_fiftytwo}
\uses{host:rank-val}
For $\eps\in(0,1/52)$ there is $x_0$ such that every rank datum $R$ with
values for a scale $s$ with $\eps(s)=\eps$ and $x\ge x_0$ has $t(R)\le
x^{400}$.
\end{lemma}

\begin{definition}[admissible rank datum]\label{host:admitted}
\lean{UpperTheory.Top.admitted}
\uses{host:host-full, host:field-rows}
% uses-lean: HostParam.Seam.Choices
A rank datum $R$ with a host with axioms $H_0$ for its frame is
\emph{admitted} if for every bound $b$ and every choice of auxiliary data
$c$ (level primes, $g_i$, $p_{\delta}$, $\mathfrak c$), field rows with bound
$b$ and $q,d>4tr+1$ imply the existence of a host $D$ for $R$ extending
$H_0$ with splitting bound $b$ and auxiliary data $c$.
\end{definition}

\begin{lemma}[rank data for hosts]\label{host:adm-classes}
\lean{UpperTheory.Top.admissibleClassesTwistedChar, UpperTheory.Top.admissible_ranks_of_CMDatum}
\uses{host:admitted, host:rank-val}
% uses-lean: GlobalTheory.EX.exists_ranks_rankVal_frame
Every rank datum with a host with axioms is admitted. For $\eps\in(0,1/52)$
there is $x_0$ such that for every scale $s$ with $\eps(s)=\eps$ and
$x\ge x_0$, every frame $\phi$ with $x^{2+8\eps}\le\mathrm{ord}\le
2x^{2+8\eps}$ and every host with axioms $H_0$ for $\phi$, there is a rank
datum with frame $\phi$ and values for $s$, which is admitted.
\end{lemma}

\begin{definition}[short norms and small primes]\label{host:short-norm}
\lean{UpperTheory.T3.ShortNorm, UpperTheory.T3.SmallOne, UpperTheory.T3.shortNorm, UpperTheory.T3.smallOne}
\texttt{SmallOne} states: there is $A\ge1$ such that for every $M\ge2$ some
prime $p\equiv1\pmod M$ has $p\le M^A$ (a Linnik-type bound).
\texttt{ShortNorm} is a lower bound, in short intervals $[X,X(1+h)]$, for
the number of primes that are norms of suitable elements of an imaginary
quadratic order of discriminant attached to a prime $d\equiv3\pmod4$ and a
modulus $M$. Both are proved in the archive (\texttt{shortNorm},
\texttt{smallOne}).
\end{definition}

\begin{lemma}[completing a host]\label{host:host-of-field}
\lean{UpperTheory.Top.hostOfField, UpperTheory.T3.hostOfFieldOf}
\uses{host:host-good, host:pair-good, host:field-rows, host:admitted, host:short-norm, host:rank-val}
% uses-lean: Formalism.HC.Frame.Shares
For all $\eps\in(0,1/52)$ and $C_0\ge1$ there are $C_{T2}$, $C_8\ge C_0$,
$c_6>0$, $C_b>0$ and $x_0$ such that the following holds. Let $s$ be a
scale with $\eps(s)=\eps$, $x\ge x_0$, $R$ a rank datum with values for $s$,
$H_0$ a host with axioms for its frame, $0<M\le C_0$ an integer and $Y$
real, with $J\le C_0$, $o\le C_0$, $Q_S\le C_0n$,
$Q_S\mid\operatorname{lcm}(2^{e+2},M)$, $Y$ a pair window for $s$, $(q,d)$
good for $(M,\eps,Y,x)$, field rows with bound $\lfloor Y^{\eps/5}\rfloor$,
and $(R,H_0)$ admitted. Then there are a host $H$ for $R$ extending $H_0$
and an integer $q_0$ with $C_{T2}(H)=C_{T2}$,
$\mathrm{HostGood}(s,C_8,4,c_6,C_b,H,q_0)$, and $TT(H)$ having shares up
to $x^{-1+\eps}$.
\end{lemma}
\begin{proof}
\texttt{hostOfField} is \texttt{hostOfFieldOf} applied to \texttt{shortNorm}
and \texttt{smallOne}. The proof is in
\arch{Development/AutomorphicArithmetic/ShortNorms/hostOfFieldOf.lean} and
was not read here.
\end{proof}

\begin{theorem}[supply from admissible classes]\label{host:supply-of-adm-of}
\lean{UpperTheory.T3.supplyOfAdmOf}
\uses{host:field-chi, host:win-pairs, host:pair-ex, host:pair-hyp, host:bounded-tw, host:ischifield-of-kerchi, host:rank-le, host:host-exw, host:gchi, host:host-of-field}
% uses-lean: UpperTheory.Top.VWinAt, Formalism.HC.EAtH, UpperTheory.Top.VVle.mono, Formalism.HC.Level.cyc
\uses{an:level-cyc, sha:vwin} % resolved
Assume the conclusion of Lemma~\ref{host:host-of-field}. Let $\psi$ be a
primitive character with \texttt{VWinAt}$(\psi,401)$, and assume: for every
$\eps\in(0,1/52)$ there is $x_0$ such that for every scale $s$ with
$\eps(s)=\eps$, $x\ge x_0$, every frame $\phi$ with
$x^{2+8\eps}\le\mathrm{ord}\le2x^{2+8\eps}$ and $o(\phi)=|\langle\psi
\rangle|$, and every host with axioms $H_0$ for $\phi$ with the kernel
condition for $\psi$, there is an admitted rank datum with frame $\phi$ and
values for $s$. Then $\mathrm{HostExW}(3/2,\mathrm{GChi}(\psi))$ holds, with
$C_5=4$.
\end{theorem}
\begin{proof}
\uses{host:field-chi, host:win-pairs, host:pair-ex, host:pair-hyp, host:bounded-tw, host:ischifield-of-kerchi, host:rank-le, host:host-of-field}
Fix $\eps$. Theorem~\ref{host:field-chi} gives $c,e_0$, and
Lemma~\ref{host:win-pairs} with Theorem~\ref{host:pair-ex} gives $J$ and
windows of exponent $3/2$. Take $C_0=J+|\langle\psi\rangle|+8N_{\psi}$ and
apply Lemma~\ref{host:host-of-field} for $C_0$; collect the thresholds of
the admissibility hypothesis and of Lemma~\ref{host:rank-le}. Choose
$A$ with $A^{\eps/5}=K$, where $K\ge\max(5,4N_{\psi})$, and let $x_1$ exceed
all thresholds, $A^J$ and $c2^{e_0}$. The system of windows is $J$ with
$W=\{x:\text{a window datum exists at }x\}$.

Let $s$ be a member scale with $x\ge x_1$ and take its window datum
$(Y,q,d,e)$. Then $Y\ge x^{1/J}\ge A$, so $Y^{\eps/5}\ge K$. By
Lemma~\ref{host:pair-hyp}, $(q,d)$ satisfies \texttt{HostHyp} and small
primes split in $K$; $q\ne3$ because $\log q>Y^{\eps/5}/2\ge5/2>\log3$;
and $e\ge e_0$ because $c2^e\ge x^{2+8\eps}\ge x>c2^{e_0}$.
Theorem~\ref{host:field-chi} with $b=\lfloor Y^{\eps/5}\rfloor\ge4N_{\psi}$
gives $\phi$ and $H_0$, the hypothesis gives an admitted $R$ with values
for $s$, and Lemma~\ref{host:host-of-field} (with $M=N_{\psi}$; the bound
$Q_S\le C_0n$ follows from $Q_S\le4N_{\psi} c2^e\le8N_{\psi} x^{2+8\eps}$)
gives the host $H$, the prime $q_0$, \texttt{HostGood} and the shares.

The explicit-formula condition \texttt{EAtH}: the index $o$, the kept
characters and the exceptional character come from
Theorem~\ref{host:field-chi}, the shares from
Lemma~\ref{host:host-of-field}, and the pointwise and log-derivative
clauses from Lemma~\ref{host:bounded-tw}, whose hypothesis on the prime
divisors of $Q_S$ holds because $Q_S\mid\operatorname{lcm}(2^{e+2},
N_{\psi})$. The Sha bound: by Lemma~\ref{host:ischifield-of-kerchi} the field
$E$ of $H$ is a field of $\psi$ at $(e,q,d)$, so the window datum gives
\texttt{VVle} with $\lfloor x^{400}\rfloor$ twists and budget
$\sqrt{\exp(x^{\eps})}=\sqrt b$, and Lemma~\ref{host:rank-le} gives
$t(R)\le x^{400}$.
\end{proof}

\begin{theorem}[supply]\label{host:supply-of-adm}
\lean{UpperTheory.Top.supplyOfAdm}
\uses{host:supply-of-adm-of, host:host-of-field}
Under the admissibility hypothesis of Theorem~\ref{host:supply-of-adm-of},
\texttt{VWinAt}$(\psi,401)$ implies $\mathrm{HostExW}(3/2,\mathrm{GChi}(\psi))$.
\end{theorem}
\begin{proof}
Theorem~\ref{host:supply-of-adm-of} applied to Lemma~\ref{host:host-of-field}.
\end{proof}

\begin{theorem}[supply for every character]\label{host:supply-chi}
\lean{UpperTheory.Top.supplyChi}
\uses{host:supply-of-adm, host:adm-classes, host:gchi, host:host-exw}
% uses-lean: UpperTheory.Top.VWinAt
\uses{sha:vwin} % resolved
For every primitive character $\psi$, \texttt{VWinAt}$(\psi,401)$ implies
$\mathrm{HostExW}(3/2,\mathrm{GChi}(\psi))$.
\end{theorem}
\begin{proof}
\uses{host:supply-of-adm, host:adm-classes}
Apply Theorem~\ref{host:supply-of-adm}; its admissibility hypothesis is
Lemma~\ref{host:adm-classes}, which does not need the conditions on $o$ or
the kernel.
\end{proof}

\begin{remark}[what the supply theorem assumes]\label{host:rem-supply}
Apart from \texttt{VWinAt}, every input of Theorem~\ref{host:supply-chi}
is a theorem of the archive: the host construction, \texttt{QH.pairEx},
\texttt{hostOfField} and the admissibility lemma. Whether those theorems
depend on further hypotheses through their unread proofs was not checked
here.
\end{remark}
